% appendix_g.tex - Detailed Baseline Reliability Analysis
\section{Detailed Baseline Reliability Analysis}\label{app:baseline-detail}

This appendix keeps the per-dataset baseline tables used by
\S\ref{sec:baseline-reliability}.  The main analysis treats the
direct-vs-recursive pattern as a fixed-budget diagnostic, not as a causal
intermittency result.  Three reproducibility checks matter most:
(i) the earlier fixed-budget protocol ranking differs across datasets,
(ii) the matched-panel aggregate-\WAPE{} view makes M5 and Rohlik more
conservative than the earlier per-series view and gives a mixed Favorita
protocol ranking, and (iii) the 200-trial M5
\LGBM{} sensitivity changes the tuning-budget interpretation without removing
the Chronos-2 M5 advantage.  The Favorita \texttt{direct\_scaled} run
should not be used as inferential evidence.

\subsection{M5 Results}

Table~\ref{tab:t53} reports the rolling-origin M5 diagnostics.  \WRMSSE{}
ranks direct LightGBM ahead of recursive LightGBM and Seasonal Naive within
this protocol.  These rows are not compared with the official M5 competition
winner because the competition used a fixed 28-day task and a different
evaluation setup.

\begin{table}[htbp]
\centering
\caption{Seasonal Naive vs LightGBM recursive vs LightGBM direct on the M5
tail-eval window.  Bold = best per horizon.}
\label{tab:t53}
\small
\begin{tabular}{llrrrrrr}
\toprule
Model & $h$ & MAE & sMAPE & \WAPE{} & \textbf{\WRMSSE{}} & Runtime (s) & Cost (USD) \\
\midrule
Seasonal Naive &  7 & 1.1590 &  75.44 & 1.3072 & 0.7758 &  236.9 & 0.025 \\
Seasonal Naive & 14 & 1.1783 &  76.06 & 1.3216 & 0.7718 &  176.9 & 0.019 \\
Seasonal Naive & 28 & 1.1897 &  76.69 & 1.3285 & 0.7661 &  139.6 & 0.015 \\
LightGBM rec   &  7 & 1.0110 & 144.44 & 1.6246 & 0.6336 &  300.3 & 0.032 \\
LightGBM rec   & 14 & 1.0391 & 144.55 & 1.7292 & 0.6639 &  301.1 & 0.032 \\
LightGBM rec   & 28 & 1.0743 & 144.48 & 1.9357 & 0.6976 &  284.3 & 0.030 \\
\textbf{LightGBM dir} &  7 & \textbf{0.9682} & 145.81 & \textbf{1.3512} & \textbf{0.5598} & 1114.4 & 0.118 \\
\textbf{LightGBM dir} & 14 & \textbf{0.9896} & 146.20 & \textbf{1.4100} & \textbf{0.5821} & 2003.2 & 0.212 \\
\textbf{LightGBM dir} & 28 & \textbf{1.0117} & 146.95 & \textbf{1.5133} & \textbf{0.5992} & 3728.8 & 0.394 \\
\bottomrule
\end{tabular}
\end{table}

On M5, direct LightGBM is better than recursive LightGBM on both \WAPE{}
and \WRMSSE{} in the fixed-budget sweep.  Seasonal Naive has lower
per-series \WAPE{} than direct LightGBM at every horizon, despite worse
\WRMSSE{}, which is why M5 \WAPE{} is treated as a denominator-pathology
warning rather than as the sole ranking metric.

\subsection{Rohlik Results}

\begin{table}[htbp]
\centering
\caption{Seasonal Naive vs LightGBM recursive vs LightGBM direct on the
Rohlik~v2 tail-eval window.  Bold = best per horizon among LightGBM
variants.}
\label{tab:t54}
\small
\begin{tabular}{llrrrrr}
\toprule
Model & $h$ & MAE & sMAPE & \WAPE{} & Runtime (s) & Cost (USD) \\
\midrule
Seasonal Naive &  7 & 40.649 & 36.46 & 0.4015 &   8.5 & 0.0009 \\
Seasonal Naive & 14 & 42.236 & 37.20 & 0.4031 &   6.9 & 0.0007 \\
Seasonal Naive & 28 & 43.543 & 38.21 & 0.4116 &   5.8 & 0.0006 \\
\textbf{LightGBM rec} &  7 & \textbf{19.335} & 92.90 & \textbf{0.3654} &  29.6 & 0.0031 \\
\textbf{LightGBM rec} & 14 & 20.972 & 94.59 & \textbf{0.3953} &  30.6 & 0.0032 \\
\textbf{LightGBM rec} & 28 & 22.815 & 97.18 & \textbf{0.4321} &  29.1 & 0.0031 \\
LightGBM dir &  7 & 19.549 & 93.34 & 0.3816 & 141.1 & 0.0149 \\
LightGBM dir & 14 & \textbf{20.210} & 94.17 & 0.4072 & 265.6 & 0.0280 \\
LightGBM dir & 28 & \textbf{21.087} & 95.01 & 0.4442 & 524.3 & 0.0553 \\
\bottomrule
\end{tabular}
\end{table}

Rohlik reverses the M5 fixed-budget \WAPE{} ranking: recursive LightGBM is
lower than direct at all three horizons, while direct is several times more
expensive.  This is a protocol-and-dataset interaction hypothesis, not a
standalone intermittency proof.

\subsection{Favorita Results And \texttt{direct\_scaled} Flag}

\begin{table}[htbp]
\centering
\caption{Seasonal Naive vs LightGBM recursive vs LightGBM direct on the
Favorita top-30k tail-eval window.  Bold = best per horizon among LightGBM
variants.}
\label{tab:t55}
\small
\begin{tabular}{llrrrrr}
\toprule
Model & $h$ & MAE & sMAPE & \WAPE{} & Runtime (s) & Cost (USD) \\
\midrule
Seasonal Naive &  7 & 5.248 & 60.76 & 0.6363 &     65 & 0.0068 \\
Seasonal Naive & 14 & 5.400 & 61.49 & 0.6473 &     49 & 0.0052 \\
Seasonal Naive & 28 & 5.644 & 62.49 & 0.6643 &     39 & 0.0041 \\
\textbf{LightGBM rec} &  7 & \textbf{2.507} & 90.92 & \textbf{0.5295} &   236 & 0.0249 \\
\textbf{LightGBM rec} & 14 & \textbf{2.574} & 91.32 & \textbf{0.5311} &   230 & 0.0242 \\
\textbf{LightGBM rec} & 28 & 2.692 & 92.13 & \textbf{0.5377} &   230 & 0.0243 \\
LightGBM dir &  7 & 2.533 & 91.38 & 0.5513 & 1{,}015 & 0.1071 \\
LightGBM dir & 14 & 2.619 & 92.17 & 0.5711 & 1{,}907 & 0.2012 \\
LightGBM dir & 28 & \textbf{2.688} & 92.75 & 0.5892 & 3{,}588 & 0.3787 \\
\bottomrule
\end{tabular}
\end{table}

In the fixed-budget Favorita sweep, recursive LightGBM is cheaper and lower
on \WAPE{} than direct LightGBM.  A follow-up
\texttt{direct\_scaled\_favorita.csv} run is marked as superseded because its
$h=7$ row changes \WAPE{} despite the same nominal 300-tree count as the base
direct run.  The matched-panel protocol check below replaces that run for
interpretation.

\subsection{Descriptive Source B Grid}\label{app:sourceb-grid}

The grid below reports per-series \WAPE{} means from the saved unmatched
Source~B result file.  Because \texttt{n\_series} varies within every
dataset--horizon cell, the table is provenance material and should not be read
as a fully paired model comparison.  The per-series denominator makes M5's
tail-sparse rows sensitive to intermittent-demand artifacts.

\begin{table*}[htbp]
\centering
\caption{Source~B \WAPE{} grid (per-series mean, rolling origin) across
models, datasets, and horizons.  Lower is better.  The grid is complete,
but \texttt{n\_series} differs across models; see
\texttt{source\_b\_n\_series\_matrix.csv}.}
\label{tab:t59}
\scriptsize
\begin{tabular}{llrrrrrrrrr}
\toprule
Model & Family & \multicolumn{3}{c}{Favorita} & \multicolumn{3}{c}{M5} & \multicolumn{3}{c}{Rohlik~v2} \\
\cmidrule(lr){3-5} \cmidrule(lr){6-8} \cmidrule(lr){9-11}
 & & $h{=}7$ & $h{=}14$ & $h{=}28$ & $h{=}7$ & $h{=}14$ & $h{=}28$ & $h{=}7$ & $h{=}14$ & $h{=}28$ \\
\midrule
Chronos-Bolt-Tiny & FM & 0.482 & 0.485 & 0.489 & 0.897 & 0.902 & 0.911 & 0.334 & 0.345 & 0.361 \\
Moirai~2.0-Small & FM & 0.477 & 0.482 & 0.487 & 0.887 & 0.893 & 0.901 & 0.324 & 0.337 & 0.354 \\
TiRex & FM & 0.525 & 0.531 & 0.539 & 0.903 & 0.907 & 0.914 & 0.323 & 0.337 & 0.353 \\
Chronos-2 & FM & 0.477 & 0.481 & 0.486 & 0.879 & 0.887 & 0.899 & 0.324 & 0.339 & 0.357 \\
TimesFM~2.5 & FM & 0.528 & 0.534 & 0.542 & 0.944 & 0.944 & 0.948 & 0.324 & 0.338 & 0.353 \\
\texttt{lightgbm\_cov} & ML\_TREE & 0.530 & 0.531 & 0.538 & 1.625 & 1.729 & 1.936 & 0.365 & 0.395 & 0.432 \\
\texttt{lightgbm\_direct} & ML\_TREE & 0.551 & 0.571 & 0.589 & 1.351 & 1.410 & 1.513 & 0.382 & 0.407 & 0.444 \\
\texttt{seasonal\_naive} & STATS & 0.636 & 0.647 & 0.664 & 1.307 & 1.322 & 1.329 & 0.402 & 0.403 & 0.412 \\
\bottomrule
\end{tabular}
\end{table*}

\subsection{Earlier Fixed-Budget Protocol Summary}

\begin{table}[htbp]
\centering
\caption{Earlier full-workload summary of the fixed-budget
direct-vs-recursive LightGBM gap using per-series \WAPE{} means
(positive = direct lower, negative = recursive lower).}
\label{tab:t58-appendix}
\small
\begin{tabularx}{\textwidth}{llrrrX}
\toprule
Dataset & Series-day zero frac. & $h = 7$ & $h = 14$ & $h = 28$ & Interpretation \\
\midrule
M5 & ${\sim}70\,\%$ & $+16.8\,\%$ & $+18.5\,\%$ & $+21.8\,\%$ & Direct lower in this protocol \\
Rohlik & ${\sim}1.2\,\%$ & $-4.4\,\%$ & $-3.0\,\%$ & $-2.8\,\%$ & Recursive lower and cheaper \\
Favorita & ${\sim}15$--$25\,\%$ & $-4.1\,\%$ & $-7.5\,\%$ & $-9.6\,\%$ & Recursive lower in fixed-budget run \\
\bottomrule
\end{tabularx}
\end{table}

\subsection{Matched-Panel Protocol Check}

The matched-panel analysis recomputes the M5, Rohlik, and Favorita LightGBM protocol
contrast on the same 100 top-volume series used for the FM-vs-baseline
paired panel.  Using aggregate \WAPE{}, direct is slightly worse than
recursive on M5 at $h=7$ and $h=14$, slightly better at $h=28$, and worse
than recursive on Rohlik at all horizons.  On Favorita, recursive is lower
at $h=7$ and $h=14$, while direct is lower at $h=28$.  The horizon-scaled
direct run equals the unscaled direct run at $h=7$ on all three datasets,
as expected when both use 300 trees.  This resolves the same-tree-count
behaviour inside the new matched-panel protocol; it does not retroactively
explain the earlier unmatched \texttt{direct\_scaled\_favorita.csv} anomaly.

The combined evidence supports a cautious hypothesis: the best LightGBM
protocol can change with the dataset, metric aggregation, sampled workload,
and training budget.  It does not support a hard intermittency threshold,
nor a general rule that direct or recursive is always preferable.

\subsection{Cost Consistency Audit}\label{app:cost-envelope}

The Azure cost tables are retained for audit, but this paper does not use them
for universal cost conclusions.  The detail tables and summary table do
not yet reconcile to one ledger.

\begin{table}[htbp]
\centering
\caption{Baseline cost consistency audit.  Differences are too large to
treat as rounding; a single run ledger is required before cost claims are
promoted from descriptive to inferential.}
\label{tab:t56}
\small
\begin{tabularx}{\textwidth}{lrrrX}
\toprule
Scope & Detail tables (USD) & Summary table (USD) & Difference & Status \\
\midrule
M5 total & 0.877 & 0.810 & $-0.067$ & mismatch requires single ledger \\
Rohlik total & 0.110 & 0.120 & $+0.010$ & mismatch requires single ledger \\
Favorita total & 0.777 & 1.119 & $+0.343$ & mismatch requires single ledger \\
Favorita direct & 0.687 & 1.030 & $+0.343$ & mismatch requires single ledger \\
M5 seasonal naive & 0.059 & 0.146 & $+0.087$ & mismatch requires single ledger \\
M5 direct & 0.724 & 0.571 & $-0.153$ & mismatch requires single ledger \\
\bottomrule
\end{tabularx}
\end{table}


The cost plots in \S\ref{sec:cost-accuracy} therefore remain cost-error
scatterplots.  The matched-panel analysis writes a local run ledger for
the completed M5, Rohlik, and Favorita reruns.  The older Azure
summary/detail tables still need a single reconciled ledger before cost
claims are promoted beyond descriptive accounting.

\subsection{Reproducibility Detail}\label{app:reproducibility}

Every baseline result should be tied to a git commit SHA, an Azure ML
pipeline run name, a registered data asset version, and a versioned
environment.  The matched-panel runs store shared series identifiers,
forecast origins, actuals, per-model predictions, and paired bootstrap
contrasts for completed M5, Rohlik, and Favorita cells.  The same requirement
applies to future GPU/larger-workload reruns and covariate-aware FM
variants before Source~B can enter a broader formal paired analysis.
